\subsection{Formal Framework: A Homeostatic Dynamical System}\label{sec:framework}

We propose that the dynamics observed across our experiments follow a common mathematical structure: a \textbf{homeostatic dynamical system} where an internal quantity $Q$ is actively regulated toward a set point $Q^*$.

\paragraph{The Governing Equation.}
Let $Q(t)$ denote the controlled quantity at training step $t$. Under homeostatic regulation, the dynamics follow:
\begin{equation}
\frac{dQ}{dt} = -k(Q - Q^*) + \eta(t)
\label{eq:homeostasis}
\end{equation}
where $k > 0$ is the relaxation rate, $Q^*$ is the equilibrium set point, and $\eta(t)$ represents perturbations from task pressure or noise. At equilibrium, $dQ/dt = 0$ implies $Q \to Q^*$ exponentially.

\paragraph{Mapping to Experiments.}
This equation unifies our observations:

\begin{itemize}
    \item \textbf{Transformer Crystallization (Section~\ref{sec:anomaly}):} $Q = \text{EDI}$, the Entropy Dispersion Index. We observe $Q^* = 0.611991643906$ across five seeds. The sharp transition ($Q \approx 0.42$ pre-grokking, then $Q^*$) represents crossing a phase boundary.
    
    \item \textbf{Le Chatelier Compensation (Section~\ref{sec:mechanism}):} When Head 0 is perturbed ($\eta \neq 0$), other heads adjust their contribution to restore EDI. The compensation score ($0.13 \pm 0.03$) quantifies $k$, the strength of the restoring force.
    
    \item \textbf{Forced Attractors:} Under dual-timescale training with explicit homeostatic loss, the system finds a new equilibrium depending on timescale separation. Testing EDI targets across 0.45–0.65, we find that high targets fail to track (the system converges to $\sim$0.46 regardless of target). However, timescale ablation reveals the attractor is \textit{timescale-dependent}:
    
    \begin{center}
    \begin{tabular}{lcc}
    \toprule
    Slow LR & Mean $Q^*$ (target 0.65) & Status \\
    \midrule
    0.01 & 0.41 & Fails to track \\
    0.05 & 0.43 & Fails to track \\
    0.10 & 0.44 & Fails to track \\
    0.50 & 0.56 & \textbf{Tracks target} \\
    \bottomrule
    \end{tabular}
    \end{center}
    
    \noindent Critically, at slow\_lr = 0.50, individual seeds reach $Q^* = 0.61$: the natural equilibrium. The ``wall'' is not architectural: with appropriate controller tuning, the system can access the full range of equilibria. This transforms the forced attractor from a fundamental constraint into an engineering parameter.
\end{itemize}

\paragraph{Predictions.}
The model makes testable predictions:
\begin{enumerate}
    \item \textbf{Exponential relaxation:} After perturbation, recovery should follow $Q(t) - Q^* \propto e^{-kt}$. (Prediction; not yet tested quantitatively. This represents a straightforward validation target for future work with minimal compute overhead.)
    \item \textbf{Equilibrium universality:} Different random seeds (same architecture/task) converge to identical $Q^*$ (confirmed to 12 decimals).
    \item \textbf{Geometric constraints:} The set of achievable $Q^*$ is bounded by architecture, not just optimizer settings.
\end{enumerate}

\paragraph{Connection to Thermodynamics.}
Equation~\ref{eq:homeostasis} is equivalent to Newton's law of cooling, where a system exchanges ``heat'' with a ``reservoir'' at temperature $Q^*$. In neural networks, $Q$ represents information-geometric state (entropy, effective rank), and $Q^*$ emerges from the interaction of task requirements and architectural capacity. The Le Chatelier response (perturbations trigger compensatory changes) is precisely how physical systems maintain equilibrium under external forcing.
